% Pursuing Stacks (cotes 134-2 to 134-8): archive cote 134-3, site page 62 (PDF page 63), typescript page 62.
% typescript. Draft compared: ps2.tex:3169-3264.
% AI-transcribed by Claude Opus 5.5 (claude-opus-5-5, adaptive thinking, effort high), Message Batches API, batch ps1 (msgbatch_01Rv5hCaNTtUt4JHx1yqoHPL), 2026-11-12; not yet checked by eye.
\scanpage{134-3/62}
\mspage{62}

(i)   For any C in (Cat)  the functor
(* * *)   $A_{/C}$ \hand{$\longrightarrow$} AxC   deduced from (*)  and $A_{/C} \hand{$\xrightarrow{\text{can}}$} A$
\xout{is}   is aspheric.
\deleted{\xout{(ii)}}\xout{For any C in (Cat), the functor (xxx) is a weak equivalence}
(ii\deleted{i})  For any C in (Cat) with final element and any a in A,  $a\mathrm{x}i^{*}(C)$ is
an aspheric element in A\^{}, i.e. for any such C, $i^{\hand{*}}(C)$ is aspheric
over the final element e \inserted{of $A^\wedge$}.
\margin{Look at \unsure{refined} \unsure{on} 63--65 vol.\ 2}
(i\over{\unsure{v}}{ii})   The element  $i^{\hand{*}}(\hand{$\Delta$}_1)$ of A\^{} is aspheric over the final object e .

\xout{We get at once the}
\hand{\uline{Remark 1.}}
Of course \unsure{\xout{a)}}the conditions in a) imply that A is aspheric (take \xout{in} C
the final category),  and hence  by an easy \inserted[typed]{\deleted{homotopy lemma}}\xout{variant of the homotopy lemma}\inserted{lemma (cf.\ below \S\ \unsure{Section (\S40)})}
we get that \inserted{for any C in (Cat),} AxC \hand{$\rightarrow$} C is a weak equivalence\inserted[typed]{(and even aspheric)}, and hence $A_{/C}$ \hand{$\rightarrow$} AxC is
a weak equivalence (because its compositum with the weak equivalence
AxC \hand{$\rightarrow$} C is a weak equivalence by assumption).  It is \xout{not} unlikely however
\margin{X}
that  a) \xout{(asphericity of A}  impl\over{\unsure{y}}{\unsure{ies}} the conditions of  b) ,  namely asphe-%
ricity (not merely weak equivalence) of (***) in  b (i).   But the opposi-%
te implication,  namely  b (i)  +  asphericity of A   implies a (i) ,
is trivial,  b\over{y}{e}cause AxC \hand{$\rightarrow$} C is aspheric.

\inserted{Proof of Theorem:}
We stated a) and b) in a way to get a visibly decreasing cascade of
\over{condi}{condi}tions; \inserted{and} moreover that the weakest in a) implies the strongest,  or
that b (ii) implies b (i) , is an immediate consequence of the standar\over{\unsure{t}}{d}
asphericity criterion for a functor between categories \inserted{\deleted{(p.38)}}\inserted{(section 30).} The only point
which is a little less formal is that  b (iii)  implies b (ii).  But
using the final object in C,  we get \inserted{(unique)} a homotopy  in (Cat) \inserted{(relative to $\Delta_1$)}
\[ \hand{$\Delta$}_1 \mathrm{x}\, C \rightarrow C , \]
between the identity map of C and the constant map with value $e_C$  - hence
by applying $i^{*}$\hand{,} a homotopy  relative to  $i^{\hand{*}}(\hand{$\Delta$}_1)$  (viewed as an ``interval''
by taking as ``endpoints'' the arrows deduced from $\hand{$\delta$}_o,\hand{$\delta$}_1 : e_{(Cat)} = \hand{$\Delta$}_o \hand{$\rightrightarrows$} \hand{$\Delta$}_1$
by applying $i^{*}$ ),  between the identity map of $i^{*}(C)$ and a constant map
of $i^{*}(C)$,  which by the  homotopy lemma implies that $i^{*}(C)$ itself is
aspheric over $e_{\hand{$A^\wedge$}}$,  qed.
\inserted[typed]{Remark}\hand{2}:  \xout{\uline{Corollary}}  the condition  b (iii) can be stated by saying that $i^{\hand{*}}(\hand{$\Delta$}_1)$
is a \uline{homotopy interval} in A\^{}.  All which needs to be checked for this, is
that  condition  b) for homotopy intervals  (p.59) namely $e_o \hand{$\cap$} e_1 = \emptyset$  is
satisfied in A\^{},  but this follows from the corresponding \xout{condition} pro-%
perty for \hand{$\Delta$}$_1$ in Cat (as the functor $i^{*}$ is left exact) and from the fact
that $i^{*}$ transforms initial element $\emptyset_{(Cat)}$ into initial element $\emptyset_{A^\wedge}$.  This
last fact is equivalent to $i(a) \neq \emptyset$ for any a in A , which is true as
i(a) has a final element.
